\section{Matrix forms of the completed intersection}
\label{sec:matrix-forms}

The compatible intersection, its endpoint jet, and the same-history restart
have now been established.  The matrices in this section reorganize those
completed rows.  The whole-terminal bounds and their compatibility have
already been proved before any matrix is introduced.

\begin{proposition}[Mixed-index transport matrix and pressure completion]
\label{prop:formal-mixed-index-matrix}
Let
\[
 \Gamma:=\Nzero\times\Nzero^3,
 \qquad
 \kappa=(n,\alpha),
 \qquad
 \lambda=(a,\beta),
\]
and write
\[
 \lambda\preceq\kappa
 \quad\Longleftrightarrow\quad
 a\le n\ \text{ and }\ \beta\le\alpha,
 \qquad
 \kappa-\lambda:=(n-a,\alpha-\beta),
 \qquad
 \binom\kappa\lambda:=\binom na\binom\alpha\beta.
\]
Set
\[
 J_\kappa:=J_{n,\alpha},
 \qquad
 \Pi_\kappa:=\Pi_{n,\alpha},
 \qquad
 \mathbf J:=(J_\lambda)_{\lambda\in\Gamma}.
\]
The block operator
\begin{equation}
 \mathbb L(J)_{\kappa,\lambda}
 :=
 \begin{cases}
 \displaystyle
 \binom\kappa\lambda
 \Pp\bigl[(J_{\kappa-\lambda}\cdot\nabla)\,\mathord\cdot\,\bigr],
 &\lambda\preceq\kappa,\\[6pt]
 0,&\lambda\npreceq\kappa,
 \end{cases}
 \label{eq:formal-mixed-transport-matrix}
\end{equation}
has the finite block-row action
\begin{equation}
 [\mathbb L(J)\mathbf J]_\kappa
 =\sum_{\lambda\preceq\kappa}\binom\kappa\lambda
   \Pp\bigl[(J_\lambda\cdot\nabla)J_{\kappa-\lambda}\bigr].
 \label{eq:formal-mixed-transport-row}
\end{equation}
Consequently the complete mixed row is the pair
\begin{align}
 J_{n+1,\alpha}+[\mathbb L(J)\mathbf J]_\kappa
 &=\nu\Delta J_\kappa,
 \label{eq:formal-mixed-projected-matrix-row}\\
 \Pi_\kappa
 &=R_iR_j\sum_{\lambda\preceq\kappa}\binom\kappa\lambda
   J_{\lambda,i}J_{\kappa-\lambda,j},
 \label{eq:formal-mixed-matrix-pressure-row}
\end{align}
while the unprojected velocity identity retains
\(\nabla\Pi_\kappa\) exactly as in
\eqref{eq:formal-mixed-velocity-row}.

Order \(\Gamma\) by increasing \(n+|\alpha|\), with any fixed ordering among
indices of equal depth.  Then \(\mathbb L(J)\) is block lower triangular.
\begin{samepage}
For
\[
 \kappa!:=n!\,\alpha!,
 \qquad
 \alpha!:=\alpha_1!\alpha_2!\alpha_3!,
 \qquad
 \widehat J_\kappa:=\frac{J_\kappa}{\kappa!},
 \qquad
 \widehat\Pi_\kappa:=\frac{\Pi_\kappa}{\kappa!},
\]
let \(D\) be the diagonal matrix with \(D_{\kappa,\kappa}=\kappa!\).  On
every finite downward-closed truncation of \(\Gamma\),
\begin{equation}
 D^{-1}\mathbb L(D\widehat J)D
 =\widehat{\mathbb L}(\widehat J),
 \qquad
 \widehat{\mathbb L}(\widehat J)_{\kappa,\lambda}
 :=
 \begin{cases}
 \Pp\bigl[(\widehat J_{\kappa-\lambda}\cdot\nabla)
            \,\mathord\cdot\,\bigr],
 &\lambda\preceq\kappa,\\
 0,&\lambda\npreceq\kappa.
 \end{cases}
 \label{eq:formal-factorial-matrix-conjugacy}
\end{equation}
\end{samepage}
Thus the factorial-normalized pressure-complete recurrence is
\begin{align}
 (n+1)\widehat J_{n+1,\alpha}
 +\sum_{\lambda\preceq\kappa}
   (\widehat J_\lambda\cdot\nabla)
   \widehat J_{\kappa-\lambda}
 +\nabla\widehat\Pi_\kappa
 &=\nu\Delta\widehat J_\kappa,
 \label{eq:formal-factorial-velocity-convolution}\\
 \widehat\Pi_\kappa
 &=R_iR_j\sum_{\lambda\preceq\kappa}
   \widehat J_{\lambda,i}\widehat J_{\kappa-\lambda,j}.
 \label{eq:formal-factorial-pressure-convolution}
\end{align}
The normalized transport blocks depend on the difference
\(\kappa-\lambda\); hence \eqref{eq:formal-factorial-matrix-conjugacy} is the
block lower-triangular convolution form of the mixed transport array.
\end{proposition}

\begin{proof}
For \(\lambda\preceq\kappa\), change the summation index from \(\lambda\)
to \(\kappa-\lambda\) and use
\(\binom\kappa{\kappa-\lambda}=\binom\kappa\lambda\).  This gives
\eqref{eq:formal-mixed-transport-row}.  Apply \(\Pp\) to
\eqref{eq:formal-mixed-velocity-row}.  The identities
\[
 \Pp J_{n+1,\alpha}=J_{n+1,\alpha},
 \qquad
 \Pp\Delta J_\kappa=\Delta J_\kappa,
 \qquad
 \Pp\nabla\Pi_\kappa=0
\]
follow from \eqref{eq:formal-mixed-divergence-row} and the Fourier symbols.
They prove \eqref{eq:formal-mixed-projected-matrix-row}; equation
\eqref{eq:formal-mixed-matrix-pressure-row} is
\eqref{eq:formal-mixed-pressure-row} in mixed-index notation.  The pressure
formula also shows that projection of the velocity row loses no pressure
coordinate.

If \(\lambda\preceq\kappa\), then
\(a+|\beta|\le n+|\alpha|\), which proves lower triangularity in the stated
ordering.  Finally,
\[
 \binom\kappa\lambda\,\lambda!\,(\kappa-\lambda)!=\kappa!.
\]
Dividing the mixed velocity and pressure rows by \(\kappa!\) proves
\eqref{eq:formal-factorial-velocity-convolution} and
\eqref{eq:formal-factorial-pressure-convolution}.  The same factorial
cancellation in each operator block proves
\eqref{eq:formal-factorial-matrix-conjugacy}.  Every row has
\((n+1)\prod_{j=1}^3(\alpha_j+1)\) supported blocks, so all displayed row
sums are finite.
\end{proof}

\begin{proposition}[Sobolev rectangle matrix and adjacent stencil]
\label{prop:sobolev-rectangle-matrix}
For \(0<T<T_*\), \(N,M\in\Nzero\), \(0\le n\le N\), and
\(m\in\mathbb Z\), set
\[
 G_{n,m}(T):=\norm{V_n}_{L^2(0,T;H^m)},
\]
and set
\[
 P_{n,M}(T):=\norm{Q_n}_{L^1(0,T;H^M)}
 +\norm{\nabla Q_n}_{L^1(0,T;H^{M-1})}.
\]
The object matrix of the finite pressure-complete rectangle is
\begin{equation}
 \mathbf R_{N,M}[u_0;T]
 :=
 \begin{pmatrix}
  V_0|_{(0,T)}&V_1|_{(0,T)}&Q_0|_{(0,T)}\\
  V_1|_{(0,T)}&V_2|_{(0,T)}&Q_1|_{(0,T)}\\
  \vdots&\vdots&\vdots\\
  V_N|_{(0,T)}&V_{N+1}|_{(0,T)}&Q_N|_{(0,T)}
 \end{pmatrix}.
 \label{eq:finite-rectangle-object-matrix}
\end{equation}
The corresponding matrix of the three column norms is
\begin{equation}
 \mathbf C_{N,M}(T)
 :=
 \begin{pmatrix}
  G_{0,M+1}(T)&G_{1,M-1}(T)&P_{0,M}(T)\\
  G_{1,M+1}(T)&G_{2,M-1}(T)&P_{1,M}(T)\\
  \vdots&\vdots&\vdots\\
  G_{N,M+1}(T)&G_{N+1,M-1}(T)&P_{N,M}(T)
 \end{pmatrix}.
 \label{eq:sobolev-pressure-complete-rectangle-matrix}
\end{equation}
If \(\mathbf G_{N,M}(T)\) denotes its first two columns, then
\begin{equation}
 \norm{\mathbf G_{N,M}(T)}_{\mathrm F}^2
 =\mathfrak R_{N,M}(u;T),
 \qquad
 \sum_{n=0}^NP_{n,M}(T)=\mathfrak Q_{N,M}(p;T).
 \label{eq:rectangle-matrix-norm-identities}
\end{equation}
For each \(0\le n\le N\), the first two entries in row \(n\) form the
adjacent Hilbert-scale stencil
\begin{equation}
 \left.
 \begin{array}{c}
  (n,M+1):\quad V_n\in L^2(0,T;H^{M+1}),\\
  (n+1,M-1):\quad V_{n+1}=\partial_tV_n
                       \in L^2(0,T;H^{M-1})
 \end{array}
 \right\}
 \longrightarrow
 (n,M):\quad V_n\in C([0,T];H^M).
 \label{eq:rectangle-adjacent-stencil}
\end{equation}
For every \(m\in\Nzero\), the Sobolev matrix measures the mixed-index array
through the exact identity
\begin{equation}
 G_{n,m}(T)^2
 =\sum_{\substack{\alpha\in\Nzero^3\\|\alpha|\le m}}
   \frac{m!}{(m-|\alpha|)!\,\alpha!}
   \norm{J_{n,\alpha}}_{L^2((0,T)\times\R^3)}^2.
 \label{eq:rectangle-mixed-index-sobolev-identity}
\end{equation}
The pressure entry in row \(n\) is \(Q_n=\Pi_{n,0}\), and its spatial rows
\(\partial_x^\alpha Q_n=\Pi_{n,\alpha}\) are fixed by
\eqref{eq:formal-mixed-matrix-pressure-row}.
\end{proposition}

\begin{proof}
The object matrix is the tuple in Definition~\ref{def:finite-rectangle},
arranged by temporal row.  The two identities in
\eqref{eq:rectangle-matrix-norm-identities} are
\eqref{eq:finite-rectangle-norm} and
\eqref{eq:finite-pressure-rectangle-norm}.  Lemma~\ref{lem:adjacent-trace}
gives \eqref{eq:rectangle-adjacent-stencil}.  Finally, Plancherel and
\[
 (1+|\xi|^2)^m
 =\sum_{\substack{\alpha\in\Nzero^3\\|\alpha|\le m}}
   \frac{m!}{(m-|\alpha|)!\,\alpha!}\,\xi^{2\alpha}
\]
give \eqref{eq:rectangle-mixed-index-sobolev-identity}.  The pressure claims
follow from the definitions of \(Q_n,\Pi_{n,\alpha}\) and
\eqref{eq:formal-mixed-matrix-pressure-row}.
\end{proof}

\begin{proposition}[Whole-terminal closed estimate matrix]
\label{prop:datum-terminal-matrix-consequence}
Under the hypotheses of Theorem~\ref{thm:datum-rectangles}, fix
\(N,M\in\Nzero\) and set
\begin{align}
 a_{n,M}
 &:=\norm{V_n}_{L^2(0,T_*;H^{M+1})},\notag\\
 b_{n,M}
 &:=\norm{V_{n+1}}_{L^2(0,T_*;H^{M-1})},
 \qquad 0\le n\le N.
 \label{eq:datum-terminal-coordinate-norms}
\end{align}
Define entrywise order on nonnegative matrices by
\(A\preceq_{\mathrm{ent}}B\) when \(A_{ij}\le B_{ij}\) for every entry.
Then, for every \(0<T<T_*\),
\begin{equation}
 \mathbf G_{N,M}(T)
 \preceq_{\mathrm{ent}}
 \mathbf G_{N,M}^*
 :=
 \begin{pmatrix}
  a_{0,M}&b_{0,M}\\
  a_{1,M}&b_{1,M}\\
  \vdots&\vdots\\
  a_{N,M}&b_{N,M}
 \end{pmatrix}.
 \label{eq:datum-terminal-matrix-bound}
\end{equation}
Moreover,
\begin{align}
 \mathfrak R_{N,M}(u;T)
 &=\norm{\mathbf G_{N,M}(T)}_{\mathrm F}^2\notag\\
 &\le\norm{\mathbf G_{N,M}^*}_{\mathrm F}^2
 =\sum_{n=0}^{N}\left(a_{n,M}^2+b_{n,M}^2\right)
 =\mathfrak R_{N,M}^*[u_0,\nu,T_*]<\infty.
 \label{eq:datum-terminal-finite-sum}
\end{align}
\end{proposition}

\begin{proof}
Proposition~\ref{prop:datum-terminal-realization} makes every
\(a_{n,M}\) and \(b_{n,M}\) finite.  Restriction from \((0,T_*)\) to
\((0,T)\) decreases each nonnegative coordinate norm and gives
\eqref{eq:datum-terminal-matrix-bound}.  Taking Frobenius norms and using
\eqref{eq:rectangle-matrix-norm-identities} together with
\eqref{eq:datum-terminal-monotone-realization} proves
\eqref{eq:datum-terminal-finite-sum}.
\end{proof}
